\fsection{Ej 6 Pág 125}
\fsubsection
[\textcolor{waveBlue2}{Investiga}]
{\textcolor{waveBlue2}{\boldit{Investiga}}}
{
	Indica si las siguientes afirmaciones son verdaderas o falsas . Justifica tu respuesta.
	\begin{enumerate}[label=\textbf{\alph*)}]
		\item Es conveniente que la huella ecológica de una persona sea lo más alta posible.
		\item La huella hídrica de un producto es la cantidad de agua que contiene .
	\end{enumerate}
}

\begin{solution}
	\begin{enumerate}[label=\textbf{\alph*)}]
		\item \textbf{Falso.} Lo conveniente es que la huella ecológica sea lo más \textbf{baja} posible. Una huella ecológica alta indica que el estilo de vida de una persona requiere más recursos naturales y superficie productiva de los que el planeta puede regenerar por sí mismo, lo que resulta insostenible a largo plazo.
		\item \textbf{Falso.} La huella hídrica no es solo el agua que el producto contiene físicamente, sino el volumen total de agua dulce que se ha utilizado, contaminado o evaporado en \textbf{todas las etapas de su cadena de producción} (lo que se conoce como "agua virtual"). Por ejemplo, el agua que contiene una manzana es mínima comparada con toda la que se usó para regar el manzano y procesar la fruta.
	\end{enumerate}
\end{solution}
